\pdfbookmark[0]{Anzeigen. Mollweide: Commentationes mathematico-philologicae. 1814.}{anzeigen-mollweide-1814}
\begin{center}
\rule{.85\textwidth}{.3pt}\par
Göttingische gelehrte Anzeigen.\quad 2. Mai 1814.\par
\rule{.85\textwidth}{.3pt}
\end{center}
\hspace*{\parindent}\textit{Commentationes mathematico-philologicae tres, sistentes explicationem duorum locorum difficilium, alterius} \textsc{Virgilii}, \textit{alterius} \textsc{Platonis}, \textit{itemque examinationem duorum mensurarum praeceptorum} \textsc{Columellae}. \textit{Adjecta est epistola ad v.\,cl.} J. G. \textsc{Schneider} \textit{de excerptis geometricis} \textsc{Epaphroditi} \textit{et} \textsc{Vitruvii Rufi} \textit{scripta ab auctore harum commentationum} \textsc{Carolo Brandano Mollweide}, \textit{astron. in acad. Lipsiensi professore.} Leipzig 1813. 122 Seiten im Oktavformat, mit einer Kupfertafel.

[Die Anzeige der ersten Abhandlung ist im Band zur Astronomie der Werke von \textsc{Gauss} abgedruckt.]

Die zweite Abhandlung über eine dunkle Stelle in \textsc{Platon}s Menon war schon 1805 der hiesigen Königlichen Gesellschaft der Wissenschaften als Manuskript vorgelegt worden, %
